\subsection{EUCLIDEAN BALLS AND SPHERES}

For each integer \(p > 0\) , let \(\mathbf{U}_p\) denote the set of all pairs \((x, y)\) of points of \(\mathbf{R}^n\) such that \(d(x, y) < 1/p\) ; the inequalities (3) show that the sets \(\mathbf{U}_p\) form a fundamental system of entourages of the uniformity of \(\mathbf{R}^n\) (cf.~Chapter IX, § 2).

From this fact and from the inequality

\[
\left| d \left(\boldsymbol {x}, \boldsymbol {y}\right) - d \left(\boldsymbol {x} ^ {\prime}, \boldsymbol {y} ^ {\prime}\right) \right| \leqslant d \left(\boldsymbol {x}, \boldsymbol {x} ^ {\prime}\right) + d \left(\boldsymbol {y}, \boldsymbol {y} ^ {\prime}\right),
\]

which is a consequence of (1), we infer that \(d(x, y)\) is uniformly continuous on \(\mathbf{R}^n \times \mathbf{R}^n\) ; consequently the norm \(||x|| = d(0, x)\) is uniformly continuous on \(\mathbf{R}^n\) .

DEFINITION 1. Given a point \(\mathbf{x}_0 \in \mathbb{R}^n\) and a real number \(r > 0\) , the open (resp. closed) Euclidean ball of \(n\) dimensions with centre \(\mathbf{x}_0\) and radius \(r\) is the set of all points \(\mathbf{x} \in \mathbb{R}^n\) such that \(d(\mathbf{x}_0, \mathbf{x}) < r\) {[}resp. \(d(\mathbf{x}_0, \mathbf{x}) \leqslant r\) {]}; the Euclidean sphere of \(n - 1\) dimensions with centre \(\mathbf{x}_0\) and radius \(r\) is the set of all \(\mathbf{x} \in \mathbb{R}^n\) such that \(d(\mathbf{x}_0, \mathbf{x}) = r\) .

When there is no risk of confusion we say simply ``ball'' (resp. ``sphere'') for ``Euclidean ball'' (resp. ``Euclidean sphere''). When n = 2, a ball of two dimensions is called a disc, and a sphere of one dimension is called a circle. When n = 1, the open (resp. closed) ball with centre \(x_{0}\) and radius r is the interval \(]x_{0} - r, x_{0} + r[\) (resp. \([x_{0} - r, x_{0} + r]\) ); the sphere with centre \(x_{0}\) and radius r is the set consisting of the two end-points \(x_{0} - r, x_{0} + r\) of these intervals.

From what has been said, the balls (open or closed) with centre \(x_{0}\) (or just those with radii 1/p, where p runs through the set of integers \textgreater0) form a fundamental system of neighbourhoods of the point \(x_{0}\) .

PROPOSITION 1. Every open (resp. closed) ball of \(\mathbf{R}^n\) is an open (resp. compact) set. The closure of an open ball is the closed ball with the same centre and the same radius; the interior of a closed ball is the open ball with the same centre and the same radius.

The open (resp. closed) ball with centre \(x_0\) and radius \(r\) is the inverse image of the interval \(] - \infty, r[\) (resp. \(] - \infty, r]\) ) under the continuous function \(d(x_0, x)\) ; it is therefore open (resp. closed and bounded, hence compact). If \(d(x_0, x) = r\) , and if \(y = x_0 + t(x - x_0)\) ( \(0 < t < 1\) ) is a point of the open segment with end-points \(x_0\) and \(x\) , we have \(d(x_0, y) = tr < r\) , and \(d(x, y) = (1 - t)r\) is as small as we please; hence \(x\) lies in the closure of the open ball with centre \(x_0\) and radius \(r\) . Again, if \(z = x + t(x - x_0)\) ( \(t > 0\) ) is a point of the open ray with origin \(x\) and direction vector \(x - x_0\) , we have

\[
d (x _ {0}, z) = (1 + t) r > r,
\]

and \(d(\boldsymbol{x}, \boldsymbol{z}) = tr\) is as small as we please; hence x is not an interior point of the closed ball with centre \(x_{0}\) and radius r.

COROLLARY. Every Euclidean sphere is a compact set and is the frontier of the open and closed balls with the same centre and the same radius.

The mapping \(x \to \frac{1}{r} (x - x_0)\) transforms the sphere (resp. open ball, closed ball) with centre \(x_0\) and radius \(r\) into the sphere (resp. open ball, closed ball) with centre \(o\) and radius \(1\) ; this sphere is denoted by \(S_{n-1}\) and is called the unit sphere in \(R^n\) . Likewise, the closed ball with centre \(o\) and radius \(1\) is denoted by \(B_n\) and is called the unit ball in \(R^n\) . The topological study of a sphere of \((n - 1)\) dimensions (resp. a closed ball of \(n\) dimensions) is thus reduced to that of \(S_{n-1}\) (resp. \(B_n\) ). For the open balls, we have the following proposition:

PROPOSITION 2. Every n-dimensional open ball is homeomorphic to \(R^{n}\) .

For the mapping \(x \to \frac{x}{1 + ||x||}\) is continuous on \(\mathbf{R}^n\) and maps \(\mathbf{R}^n\) onto the open ball with centre \(0\) and radius \(1\) ; moreover, from \(y = \frac{x}{1 + ||x||}\) we deduce \(x = \frac{y}{1 - ||y||}\) , so that the mapping is bijective and bicontinuous.

Let \(\mathbf{R}_n^*\) denote the complement of o in \(\mathbf{R}^n\) .

PROPOSITION 3. The space \(\mathbf{R}_n^*\) is homeomorphic to the product of \(\mathbf{S}_{n-1}\) and the space \(\mathbf{R}_+^*\) of real numbers \(>0\) .

For every point \(x \neq 0\) can be written uniquely in the form \(tz\) , where \(t > 0\) and \(||z|| = 1\) , since \(x = tz\) implies \(t = ||x||\) and \(z = x / ||x||\) . Since \(tz\) is continuous on the product \(\mathbf{R} \times \mathbf{R}^n\) and hence a fortiori on \(\mathbf{R}_+^* \times \mathbf{S}_{n-1}\) , and since \(||x||\) and \(\frac{1}{||x||}\) are continuous on \(\mathbf{R}_n^*\) , the proposition is proved.

The mapping \(x \to x / ||x||\) is called the central projection of \(\mathbf{R}_n^*\) onto \(\mathbf{S}_{n-1}\) . One defines in the same way the central projection of the complement of a point \(a\) onto a sphere with centre \(a\) .

COROLLARY 1. The sphere \(S_{n-1}\) is homeomorphic to the quotient of \(R_{n}^{*}\) by the equivalence relation whose classes are the open rays with origin o.

These classes can also be defined as the classes of intransitivity, other than \(\{0\}\) , of the group of homotheties of ratio \textgreater0.

COROLLARY 2. The space \(\mathbf{R}_n^*\) is homeomorphic to \(\mathbf{R} \times \mathbf{S}_{n-1}\) .

For \(\mathbf{R}_{+}^{*} = ]0, +\infty[\) is homeomorphic to \(\mathbf{R}\) (Chapter IV, § 4, no. 1, Proposition 1).

Remark. These propositions are not peculiar to Euclidean balls, but can be extended to a whole category of compact neighbourhoods of o in \(R^{n}\) (see Exercise 12).

The sets \(S_{n-1}\) and \(B_{n}\) are evidently invariant under all orthogonal transformations. If V is a p-dimensional vector subspace in \(R^{n}\) , there exists an orthogonal transformation which transforms V into a p-dimensional coordinate variety; hence \(V \cap S_{n-1}\) (resp. \(V \cap B_{n}\) ) is homeomorphic to \(S_{p-1}\) (resp. \(B_{p}\) ).

